% TOP_PROBLEM: {"id": 20001036, "title": "Rainbow reachability and the multi-label core obstruction", "category_id": 2, "category": {"id": 2, "name": "combinatorics", "display_name": "Combinatorics", "description": "Counting problems, graph theory, discrete structures.", "slug": "combinatorics", "order_index": 2, "created_at": "2026-07-31T15:26:25.671Z"}, "status": "partially_solved", "problem_number": "AIM-COMBINATORICS-0161", "set_id": 14}
% FIRST_POSTED: 2026-10-01
% ENTRY_KIND: statement_only
% UnsolvedMath ID 20001036; source catalogue code AIM-COMBINATORICS-0161
% !TeX program = lualatex
% Definition and sources only; this file is not an LLM solution attempt.
\documentclass[11pt,a4paper]{article}
\usepackage[margin=25mm]{geometry}
\usepackage{fontspec}
\setmainfont{lmroman10-regular.otf}[BoldFont=lmroman10-bold.otf,
 ItalicFont=lmroman10-italic.otf,BoldItalicFont=lmroman10-bolditalic.otf]
\usepackage{amsmath,amssymb,mathtools}
\usepackage{unicode-math}
\setmathfont{latinmodern-math.otf}
\AtBeginDocument{\let\setminus\smallsetminus}
\usepackage{xurl,hyperref}
\hypersetup{colorlinks=true,urlcolor=blue}
\setcounter{secnumdepth}{0}
\setlength{\parindent}{0pt}
\setlength{\parskip}{5pt}
\setlength{\emergencystretch}{3em}
\begin{document}
\section*{Rainbow reachability and the multi-label core obstruction}\label{20001036}
{\small UnsolvedMath ID 20001036 · AIM-COMBINATORICS-0161 · Combinatorics · AIM Workshop Problem Lists}
\subsection{Definitions and mathematical statement}
Let $G=(V,E)$ be a finite simple loopless digraph with $V\ne\varnothing$,
and choose subsets $E_1,\ldots,E_k\subseteq E$ for some $k\ge1$.
These subsets need not be disjoint. An arc $e$ has the available labels
$S_e=\{i:e\in E_i\}$. A directed path or directed cycle is rainbow
if one can assign to each of its arcs a label in $S_e$ so that all
assigned labels are distinct. Labels are chosen for that occurrence
of the path or cycle; they need not increase along it.
An arc having no labels cannot be used in a rainbow structure.

Set $G_i=(V,E_i)$, and let $d_i^+(v)$ be its outdegree at $v$.
Let $R(v)$ be the set of vertices other than $v$ reachable from $v$
by a positive-length rainbow directed path. Paths have distinct
vertices, while cycles have distinct vertices apart from their common
start and end. The proposed alternative is
\[
 \text{there exists a rainbow directed cycle,}
 \quad\text{or}\quad
 \exists v\in V:\ |R(v)|\ge\sum_{i=1}^k d_i^+(v).
\]
The sum counts a multi-labelled outgoing arc once for each of its
labels; reachable vertices are counted only once. This distinction
is part of the force of the assertion. There is no restriction on
opposite arcs, and no global assignment of a single color to every
arc is assumed.
\subsection{Short English statement}
Must a directed graph with overlapping edge-label sets have either a rainbow cycle or a vertex with sufficiently large rainbow reachability?
\subsection{Sources}
\begin{itemize}
\item[S1] ULAM.AI and the UnsolvedMath Contributors, supplied problems.json, record 20001036 (AIM-COMBINATORICS-0161), AIM Workshop Problem Lists. Catalogue identity and source transcription; CC BY 4.0.
\url{https://huggingface.co/datasets/ulamai/UnsolvedMath}.
\item[S2] American Institute of Mathematics, The Caccetta-Haggkvist conjecture, item 6. Original workshop question underlying AIM-COMBINATORICS-0161.
\url{https://aimath.org/WWN/caccetta/caccetta.pdf}.
\end{itemize}
\end{document}
